\chapter*{List of Abbreviations}
\addcontentsline{toc}{chapter}{List of Abbreviations}

The following abbreviations and acronyms are used throughout this thesis. They
are collected here for reference; each is also expanded at its first occurrence
in the running text.

\begin{longtable}{L{3.2cm} L{10.5cm}}
\toprule
\textbf{Abbreviation} & \textbf{Meaning} \\
\midrule
\endfirsthead
\toprule
\textbf{Abbreviation} & \textbf{Meaning} \\
\midrule
\endhead
\midrule
\multicolumn{2}{r}{\small\itshape continued on next page} \\
\endfoot
\bottomrule
\endlastfoot

API & Application Programming Interface \\
bcrypt & Blowfish-based adaptive password-hashing function (used at cost~10) \\
CD & Continuous Deployment \\
CE & Community Edition (as in Judge0~CE, the self-hostable code sandbox) \\
CI & Continuous Integration \\
CORS & Cross-Origin Resource Sharing \\
CSRF & Cross-Site Request Forgery \\
CWE & Common Weakness Enumeration \\
DDL & Data Definition Language (SQL schema statements) \\
GIDS & Google Identity Services \\
HS256 & HMAC-SHA-256, the symmetric JWT signing algorithm used by APEX \\
HTTP & Hypertext Transfer Protocol \\
HTTPS & Hypertext Transfer Protocol Secure \\
IDOR & Insecure Direct Object Reference \\
IOI & International Olympiad in Informatics \\
JSON & JavaScript Object Notation \\
JSONB & Binary JSON column type in PostgreSQL \\
JWT & JSON Web Token \\
LMS & Learning Management System \\
LOC & Lines of Code \\
MCQ & Multiple-Choice Question \\
OIDC & OpenID Connect \\
ORM & Object-Relational Mapping (here, GORM) \\
OWASP & Open Worldwide Application Security Project \\
REST & Representational State Transfer \\
SMTP & Simple Mail Transfer Protocol \\
SPA & Single-Page Application \\
SQL & Structured Query Language \\
SUS & System Usability Scale \\
UI & User Interface \\
UX & User Experience \\
WAF & Web Application Firewall \\
XSS & Cross-Site Scripting \\

\end{longtable}
